% ==============================================================================
% appendices/appD_instructors_and_research_units.tex — ภาคผนวก ง
% ==============================================================================
\chapter{คณะวิทยากร ผู้ร่วมพัฒนาหลักสูตร และหน่วยวิจัยเฉพาะทาง}
\label{app:instructors_and_research_units}

คู่มือปฏิบัติการและการฝึกอบรมหลักสูตร \textbf{``ปัญญาประดิษฐ์ฝังตัวและเอดจ์เอไอโอทีเพื่อการเกษตรและสิ่งแวดล้อม (Embedded AI \& Edge AIoT for Agriculture and Environment)''} ได้รับการออกแบบ วิจัย และพัฒนาขึ้นโดยความร่วมมือของคณาจารย์ผู้ทรงคุณวุฒิ นักวิจัย และห้องปฏิบัติการวิจัยเฉพาะทาง คณะวิทยาศาสตร์และเทคโนโลยี มหาวิทยาลัยราชภัฏรำไพพรรณี เพื่อขับเคลื่อนการประยุกต์ใช้เทคโนโลยีปัญญาประดิษฐ์และไอโอทีในการยกระดับคุณภาพชีวิต เกษตรกรรมแม่นยำ และการอนุรักษ์สิ่งแวดล้อมอย่างยั่งยืน

\academicsection{คณะวิทยากรและผู้ร่วมพัฒนาหลักสูตร (Instructor \& Development Team)}
\label{sec:instructor_team}

\begin{enumerate}
    \item \textbf{ผู้ช่วยศาสตราจารย์ ดร.ชีวะ ทัศนา} (Asst. Prof. Dr. Chewa Thassana)\\
    \textit{บทบาทในหลักสูตร:} วิทยากรหลักและหัวหน้าโครงการ (Lead Instructor \& Project Director)\\
    \textit{ความเชี่ยวชาญ:} ปัญญาประดิษฐ์อธิบายได้ (xAI), ฟิสิกส์เกษตรดิจิทัล, สเปกโทรสโกปีดินเชิงคำนวณ, ระบบปัญญาประดิษฐ์ฝังตัวและอุปกรณ์ขอบ (Embedded \& Edge AIoT)\\
    \textit{สังกัด:} สาขาวิชาปัญญาประดิษฐ์เพื่อเกษตรดิจิทัลและสิ่งแวดล้อม คณะวิทยาศาสตร์และเทคโนโลยี มหาวิทยาลัยราชภัฏรำไพพรรณี

    \item \textbf{ผู้ช่วยศาสตราจารย์ อรรถกร คำฉัตร} (Asst. Prof. Attakorn Khamchat)\\
    \textit{บทบาทในหลักสูตร:} วิทยากรและผู้ร่วมพัฒนาหลักสูตร (Instructor \& Co-Developer)\\
    \textit{ความเชี่ยวชาญ:} วิศวกรรมระบบเกษตรอัจฉริยะ, ระบบอัตโนมัติภาคสนาม, การประมวลผลข้อมูลทางวิศวกรรมเกษตร\\
    \textit{สังกัด:} คณะวิทยาศาสตร์และเทคโนโลยี มหาวิทยาลัยราชภัฏรำไพพรรณี

    \item \textbf{ผู้ช่วยศาสตราจารย์ ดร.จิรภัทร จันทมาลี} (Asst. Prof. Dr. Jirapat Janthamalee)\\
    \textit{บทบาทในหลักสูตร:} วิทยากรและผู้ร่วมพัฒนาหลักสูตร (Instructor \& Co-Developer)\\
    \textit{ความเชี่ยวชาญ:} จุลชีววิทยาประยุกต์ทางการเกษตร, ความหลากหลายทางชีวภาพของจุลินทรีย์ดิน, โรคพืชวิทยาและปฏิสัมพันธ์ทางชีวภาพ\\
    \textit{สังกัด:} คณะวิทยาศาสตร์และเทคโนโลยี มหาวิทยาลัยราชภัฏรำไพพรรณี

    \item \textbf{ผู้ช่วยศาสตราจารย์ ดร.ณมนรัก คำฉัตร} (Asst. Prof. Dr. Namonrak Khamchat)\\
    \textit{บทบาทในหลักสูตร:} วิทยากรและผู้ร่วมพัฒนาหลักสูตร (Instructor \& Co-Developer)\\
    \textit{ความเชี่ยวชาญ:} วิทยาศาสตร์และเทคโนโลยีสิ่งแวดล้อม, การประเมินคุณภาพทรัพยากรธรรมชาติและสิ่งแวดล้อมเชิงพื้นที่, การจัดการมลพิษ\\
    \textit{สังกัด:} คณะวิทยาศาสตร์และเทคโนโลยี มหาวิทยาลัยราชภัฏรำไพพรรณี

    \item \textbf{รองศาสตราจารย์ ดร.นิภัทร เปี่ยมอรุณ} (Assoc. Prof. Dr. Nipat Peamaroon)\\
    \textit{บทบาทในหลักสูตร:} วิทยากรผู้ทรงคุณวุฒิและที่ปรึกษาโครงการ (Senior Instructor \& Project Advisor)\\
    \textit{ความเชี่ยวชาญ:} วิทยาศาสตร์ประยุกต์, เทคโนโลยีการวิเคราะห์ขั้นสูง, การออกแบบการวิจัยและการถ่ายทอดองค์ความรู้เชิงวิชาการ\\
    \textit{สังกัด:} คณะวิทยาศาสตร์และเทคโนโลยี มหาวิทยาลัยราชภัฏรำไพพรรณี

    \item \textbf{รองศาสตราจารย์ ดร.นันทพร มูลรังษี} (Assoc. Prof. Dr. Nantaporn Moonrangsee)\\
    \textit{บทบาทในหลักสูตร:} วิทยากรผู้ทรงคุณวุฒิและที่ปรึกษาโครงการ (Senior Instructor \& Project Advisor)\\
    \textit{ความเชี่ยวชาญ:} ฟิสิกส์เกษตร, พลศาสตร์สิ่งแวดล้อม, การตรวจวัดและบริหารจัดการทรัพยากรเชิงพื้นที่\\
    \textit{สังกัด:} คณะวิทยาศาสตร์และเทคโนโลยี มหาวิทยาลัยราชภัฏรำไพพรรณี
\end{enumerate}

\academicsection{หน่วยวิจัยและห้องปฏิบัติการเฉพาะทาง (Specialized Research Units)}
\label{sec:research_units}

หลักสูตรนี้ได้รับการสนับสนุนและดำเนินงานทางวิชาการร่วมกับ 2 หน่วยวิจัยเฉพาะทางของคณะวิทยาศาสตร์และเทคโนโลยี มหาวิทยาลัยราชภัฏรำไพพรรณี:

\begin{tcolorbox}[colback=teal!5!white,colframe=teal!65!black,title=\textbf{1. หน่วยวิเคราะห์คุณภาพสิ่งแวดล้อมสำหรับพื้นที่ (LEQS)},arc=2mm]
    \textbf{Laboratory for Environmental Quality for Space/Site analysis (LEQS)}\\
    หน่วยปฏิบัติการที่มุ่งเน้นการตรวจวิเคราะห์ ตรวจสอบ และติดตามตรวจสอบคุณภาพสิ่งแวดล้อม (ดิน น้ำ อากาศ และชีวมวล) เชิงพื้นที่ เพื่อสร้างฐานข้อมูลสิ่งแวดล้อมที่มีความแม่นยำสูง สำหรับการจัดการพื้นที่เกษตรปลอดภัย การประเมินผลกระทบสิ่งแวดล้อม และการสนับสนุนมาตรฐานผลผลิตทางการเกษตรในภาคตะวันออก
\end{tcolorbox}

\begin{tcolorbox}[colback=orange!5!white,colframe=orange!80!black,title=\textbf{2. หน่วยวิจัยปัญญาประดิษฐ์ฝังตัวและเอดจ์เอไอโอทีเพื่อเกษตรดิจิทัลและสิ่งแวดล้อม},arc=2mm]
    \textbf{Embedded AI \& Edge AIoT for Digital Agriculture and Environment Research Unit}\\
    หน่วยวิจัยที่มุ่งพัฒนานวัตกรรมด้านเทคโนโลยีปัญญาประดิษฐ์บนอุปกรณ์ขอบ (Edge AI), การประมวลผลข้อมูลภาพทางการเกษตร (Computer Vision / Bio-Spectral Imaging), เซนเซอร์เคมีดิน In-Situ (RS485 Modbus / NPK-pH-Moisture), และระบบโทรมาตรไร้สาย (BLE / Wi-Fi IoT Telemetry) เพื่อการวินิจฉัยสุขภาพพืชและตรวจวิเคราะห์เนื้อดินได้แบบทันทีภาคสนาม (Real-time In-field Diagnostics)
\end{tcolorbox}

\begin{center}
    \vspace{6pt}
    {\fontsize{9.5}{13}\selectfont\color{gray!70!black}
    \textbf{คณะวิทยาศาสตร์และเทคโนโลยี มหาวิทยาลัยราชภัฏรำไพพรรณี}\\
    41 หมู่ 5 ตำบลท่าช้าง อำเภอเมือง จังหวัดจันทบุรี 22000\\
    โทรศัพท์: 0-3931-9111 \enspace\textbullet\enspace เว็บไซต์: \url{https://scitech.rbru.ac.th}}
\end{center}
